\section{Rectangles of arbitrary odd height}\label{sec:odd}

Fix $m=2k+1\ge1$. Let $a_m(n)$ be the number of domino tilings
of an $m\times2n$ rectangle up to its full symmetry group,
with $a_m(0)=1$. A nonempty rectangle of odd height and even
width is never a square, so its symmetry group is still
$\{e,h,v,\rho\}$.
Use $T_m,\chi_m,\evec$ as defined in Section \ref{sec:half}.
The following is an explicit specialization of the general
recurrence framework discussed by DeFord \cite{deford}.
We include the proof to specify all symmetry boundary terms.

\begin{theorem}\label{thm:odd}
For every $m=2k+1\ge1$ and $n\ge0$,
\begin{equation}\label{eq:odd-count}
 a_m(n)=\frac14\left[
 \evec^{\trans}T_m^{2n}\evec+
 \evec^{\trans}T_k^{2n}\evec+
 \evec^{\trans}T_m^n(\one+\chi_m)\right].
\end{equation}
Equivalently,
\begin{equation}\label{eq:odd-gf}
\begin{split}
 A_m(x):=\sum_{n\ge0}a_m(n)x^n
 =\frac14\bigl[&
 \evec^{\trans}(I-xT_m^2)^{-1}\evec\\
 &+\evec^{\trans}(I-xT_k^2)^{-1}\evec\\
 &+\evec^{\trans}(I-xT_m)^{-1}(\one+\chi_m)\bigr].
\end{split}
\end{equation}
For each fixed odd height, $A_m(x)$ is rational, and its
coefficient sequence satisfies a constant-coefficient linear
recurrence.
\end{theorem}
\begin{proof}
The identity contribution is given by \eqref{eq:profile}.
Lemma \ref{lem:h-odd} gives the horizontal-axis reflection
contribution, and Lemma \ref{lem:half} gives the other two.
The same double count as in Lemma \ref{lem:burnside} proves
\eqref{eq:odd-count}.
Summing over $n$ and using
$\sum_{n\ge0}x^nM^n=(I-xM)^{-1}$ proves \eqref{eq:odd-gf}.
Every matrix is finite, and $I-xM$ has constant term $I$.
Its inverse exists formally, and its entries are rational
by the adjugate formula.
Multiplying the resulting generating function by a denominator
with nonzero constant term and comparing sufficiently high
coefficients gives the asserted recurrence.
\end{proof}

The vectors $\evec$ in the different summands belong to the
corresponding state spaces; when $k=0$, the convention
$T_0=(1)$ applies.
Lemma \ref{lem:half} itself holds for even heights as well.
However, \eqref{eq:h-odd} uses a middle row and cannot be
applied unchanged at even height.
That case requires a separate treatment of vertical dominoes
crossing the horizontal axis. For a square, quarter-turns and
diagonal reflections must also be included to count classes
under the full symmetry group.
Neither extension is asserted here.
